\documentclass[12pt,letterpaper]{article}
\usepackage{graphicx,textcomp}
\usepackage{natbib}
\usepackage{setspace}
\usepackage{fullpage}
\usepackage{color}
\usepackage[reqno]{amsmath}
\usepackage{amsthm}
\usepackage{fancyvrb}
\usepackage{amssymb,enumerate}
\usepackage[all]{xy}
\usepackage{endnotes}
\usepackage{lscape}
\newtheorem{com}{Comment}
\usepackage{float}
\usepackage{hyperref}
\newtheorem{lem} {Lemma}
\newtheorem{prop}{Proposition}
\newtheorem{thm}{Theorem}
\newtheorem{defn}{Definition}
\newtheorem{cor}{Corollary}
\newtheorem{obs}{Observation}
\usepackage[compact]{titlesec}
\usepackage{dcolumn}
\usepackage{booktabs}
\usepackage{tikz}
\usepackage{siunitx}   
\usetikzlibrary{arrows}
\usepackage{multirow}
\usepackage{xcolor}
\usepackage[table]{xcolor} 
\usepackage{tabularray}
\UseTblrLibrary{booktabs}
\newcolumntype{.}{D{.}{.}{-1}}
\newcolumntype{d}[1]{D{.}{.}{#1}}
\definecolor{light-gray}{gray}{0.65}
\usepackage{url}
\usepackage{listings}
\usepackage{color}

\definecolor{codegreen}{rgb}{0,0.6,0}
\definecolor{codegray}{rgb}{0.5,0.5,0.5}
\definecolor{codepurple}{rgb}{0.58,0,0.82}
\definecolor{backcolour}{rgb}{0.95,0.95,0.92}

\lstdefinestyle{mystyle}{
	backgroundcolor=\color{backcolour},   
	commentstyle=\color{codegreen},
	keywordstyle=\color{magenta},
	numberstyle=\tiny\color{codegray},
	stringstyle=\color{codepurple},
	basicstyle=\footnotesize,
	breakatwhitespace=false,         
	breaklines=true,                 
	captionpos=b,                    
	keepspaces=true,                 
	numbers=left,                    
	numbersep=5pt,                  
	showspaces=false,                
	showstringspaces=false,
	showtabs=false,                  
	tabsize=2
}
\lstset{style=mystyle}
\newcommand{\Sref}[1]{Section~\ref{#1}}
\newtheorem{hyp}{Hypothesis}

\title{Problem Set 3}
\date{Due: March 25, 2026}
\author{Applied Stats II}


\begin{document}
	\maketitle
	\section*{Instructions}
	\begin{itemize}
	\item Please show your work! You may lose points by simply writing in the answer. If the problem requires you to execute commands in \texttt{R}, please include the code you used to get your answers. Please also include the \texttt{.R} file that contains your code. If you are not sure if work needs to be shown for a particular problem, please ask.
\item Your homework should be submitted electronically on GitHub in \texttt{.pdf} form.
\item This problem set is due before 23:59 on Wednesday March 25, 2026. No late assignments will be accepted.
	\end{itemize}

	\vspace{.25cm}
\section*{Question 1}
\vspace{.25cm}
\noindent We are interested in how governments' management of public resources impacts economic prosperity. Our data come from \href{https://www.researchgate.net/profile/Adam_Przeworski/publication/240357392_Classifying_Political_Regimes/links/0deec532194849aefa000000/Classifying-Political-Regimes.pdf}{Alvarez, Cheibub, Limongi, and Przeworski (1996)} and is labelled \texttt{gdpChange.csv} on GitHub. The dataset covers 135 countries observed between 1950 or the year of independence or the first year forwhich data on economic growth are available ("entry year"), and 1990 or the last year for which data on economic growth are available ("exit year"). The unit of analysis is a particular country during a particular year, for a total $>$ 3,500 observations. 

\begin{itemize}
	\item
	Response variable: 
	\begin{itemize}
		\item \texttt{GDPWdiff}: Difference in GDP between year $t$ and $t-1$. Possible categories include: "positive", "negative", or "no change"
	\end{itemize}
	\item
	Explanatory variables: 
	\begin{itemize}
		\item
		\texttt{REG}: 1=Democracy; 0=Non-Democracy
		\item
		\texttt{OIL}: 1=if the average ratio of fuel exports to total exports in 1984-86 exceeded 50\%; 0= otherwise
	\end{itemize}
	
\end{itemize}
\newpage
\noindent Please answer the following questions:

\begin{enumerate}
	\item Construct and interpret an unordered multinomial logit with \texttt{GDPWdiff} as the output and "no change" as the reference category, including the estimated cutoff points and coefficients. 
	
	First step is to load in the data and create a categorical variable that transforms GPDWdiff into a variable with 3 categories - 'Positive' (if diff is greater than 0), 'Negative' (if diff is less than -0) and 'No Change' (if diff is 0).
	
		\lstinputlisting[language=R, firstline=41, lastline= 55]{PS03_EB.R}
		\begin{verbatim}
		No Change  Positive  Negative 
		16    	  2600    	  1105  
		\end{verbatim}
	Despite the baseline 'No Change' group having only 16 data points, the variable now has factors that are unordered with 'No Change' as the reference category to run an unordered multinomial logit.
	
		\lstinputlisting[language=R, firstline=61, lastline= 75]{PS03_EB.R}
			
		% Table created by stargazer v.5.2.3 by Marek Hlavac, Social Policy Institute. E-mail: marek.hlavac at gmail.com
		% Date and time: Sat, Mar 21, 2026 - 08:30:29
		\begin{table}[!htbp] \centering 
			\caption{Unordered Multinomial Logit} 
			\label{} 
			\begin{tabular}{@{\extracolsep{5pt}}lcc} 
				\\[-1.8ex]\hline 
				\hline \\[-1.8ex] 
				& \multicolumn{2}{c}{DV} \\ 
				\cline{2-3} 
				\\[-1.8ex] & Positive & Negative \\ 
				\\[-1.8ex] & (1) & (2)\\ 
				\hline \\[-1.8ex] 
				REG & 1.769$^{**}$ & 1.379$^{*}$ \\ 
				& (0.767) & (0.769) \\ 
				& & \\ 
				OIL & 4.576 & 4.784 \\ 
				& (6.885) & (6.885) \\ 
				& & \\ 
				Constant & 4.534$^{***}$ & 3.805$^{***}$ \\ 
				& (0.269) & (0.271) \\ 
				& & \\ 
				\hline \\[-1.8ex] 
				Akaike Inf. Crit. & 4,690.770 & 4,690.770 \\ 
				\hline 
				\hline \\[-1.8ex] 
				\textit{Note:}  & \multicolumn{2}{r}{$^{*}$p$<$0.1; $^{**}$p$<$0.05; $^{***}$p$<$0.01} \\ 
			\end{tabular} 
		\end{table} 
			
			
		\begin{verbatim} 
						  REG OIL No Change Positive Negative
						1   0   0    0.0072   0.6696   0.3232
						2   1   0    0.0014   0.7526   0.2460
						3   0   1    0.0001   0.6273   0.3727
						4   1   1    0.0000   0.7131   0.2869
		\end{verbatim}
			
		Above are the results from estimating an unordered multinomial logit model with 'No Change' as the reference category and 'Positive' and 'Negative' differences in GDP between the year and the year before as a function of Regime (Democracy = 1) and Oil (Average Ratio of Fuel Exports $<$ 50\% = 1).
		
		\textbf{Intercepts} \\
		The intercepts represent the baseline log odds of going from the reference category (No change) to another category when all the covariates are zero. The baseline predicted log-odd of falling into the Positive category from No change is $\beta_0 = 4.534$, and the baseline log-odds of falling into the Negative category from No Change ($\beta_0 = 3.805$). \\
		
	
		\textbf{Coefficients} \\
		Regime type (i.e. being a democracy) is associated with an increase in the log odds of both 'Positive' and 'Negative' GDP difference outcomes relative to 'No Change', but only statistically significant in the 'Positive' category.  Holding all else constant, a country being a democracy is associated with an increase of 1.769 in the the log-odds of having a 'Positive' difference in GDP versus 'No change', i.e. the odds are roughly 5.86 times higher ($e^{1.769})\approx5.86$) for democracies than non-democracies for being in the 'Positive' category than 'No Change'. This effect is statistically significant at the 5\% level ($p <0.05$).

		Holding all else constant, a country being a democracy is associated with an increase of 1.379 in the the log-odds of having a 'Negative' difference in GDP versus 'No change', i.e. the odds are roughly 3.97 times higher ($e^{1.379}\approx3.970$) for democracies than non-democracies for being in the 'Negative' category than 'No Change'. This effect however is only significant at the 10\% level and does not meet the standard criteria to be deemed statistically significant. 
		
		Intuitively democracy is associated with a greater likelihood of positive GDP change from the previous year, although there is weak evidence that being a democracy is also associated to a negative change potentially indicating democracies are more likely to have shifts in GDP year on year in either direction compared to non-democracies. \\
		 
		 
		There is no statistically significant evidence that being a dominant oil exporter has any effect on the log-odds of being in any category relevant to 'No Change'.  Oil does not reliably help predict category outcomes of GDP shift.  \\
		
		\textbf{Predicted Probabilities} \\
		As these coefficients are expressed in log-odds relative to the baseline category, their real-world importance is best evaluated using predicted probabilities. There are four possible combinations of covariates where each row displays the predicted probabilities for the each outcome.  \\
		
		Overall we can see that across all countries the most likely outcome is having a Positive GDP difference, then Negative and it is very unlikely to have No Change (probably due to how specific the zero qualification is). 
		Here  
		First the impact of regime type can be explored by comparing the first two rows and the last two where oil is constant at 0 and then 1 respectively. 
		When Oil = 0, the predicted probabilities of positive changes in GDP increases (.67 to .75) whilst the probability for a negative change decreases (.32 to .24), indicating that democracy reduces the probability of negative changes in GDP by about 8\% and increases the probability of positive change GDP outcome by about 8\%. This effect is similar for when Oil = 1, as predicted probabilities for negative GDP difference reduces by about 8\% for democracies with a corresponding increase for positive differences. Intuitively, Democracy consistently shifts the probabilities away from negative GDP difference outcomes, irrespective of whether they are a dominant oil exporter or not. \\
		
		Next, the impact of a country being a dominant oil exporter on likelihood of experiencing a difference in GDP change can be seen by comparing the odd and even rows. We can clearly see that oil increases the predicted probability of negative changes in GDP outcomes and decreases the predicted probability of positive outcomes for both democracies (negative .24 to .28, positive .75 to .71) and non-democracies  (negative .32 to .37, positive .67 to .62). Being a dominant oil exporter raises the probability of negative GDP changes by roughly 4-5\%. Intuitively we can see that oil is associated with a higher probability of negative GDP differences from previous year, with a lower probability of positive changes. 
		
	\item Construct and interpret an ordered multinomial logit with \texttt{GDPWdiff} as the outcome variable, including the estimated cutoff points and coefficients. \\
	
	Now repeat using an ordered multinomial logit. 	Firstly need to specify the factor is now ordered, and set the order Negative$<$No Change$<$Positive to specify monotonic progression and meaningful order. 
		\lstinputlisting[language=R, firstline=82, lastline= 90]{PS03_EB.R}
	Then running the ordered regression. 
		\lstinputlisting[language=R, firstline=92, lastline= 106]{PS03_EB.R}			
				% Table created by stargazer v.5.2.3 by Marek Hlavac, Social Policy Institute. E-mail: marek.hlavac at gmail.com
				% Date and time: Sat, Mar 21, 2026 - 08:51:49
				\begin{table}[!htbp] \centering 
					\caption{Ordered Multinomial Logit} 
					\label{} 
					\begin{tabular}{@{\extracolsep{5pt}}lc} 
						\\[-1.8ex]\hline 
						\hline \\[-1.8ex] 
						& \multicolumn{1}{c}{DV} \\ 
						\cline{2-2} 
						\\[-1.8ex] & GDPWdiff\_cat \\ 
						\hline \\[-1.8ex] 
						REG & 0.398$^{***}$ \\ 
						& (0.075) \\ 
						& \\ 
						OIL & $-$0.199$^{*}$ \\ 
						& (0.116) \\ 
						& \\ 
						\hline \\[-1.8ex] 
						Observations & 3,721 \\ 
						\hline 
						\hline \\[-1.8ex] 
						\textit{Note:}  & \multicolumn{1}{r}{$^{*}$p$<$0.1; $^{**}$p$<$0.05; $^{***}$p$<$0.01} \\ 
					\end{tabular} 
				\end{table} 

				
			\begin{verbatim}
				Intercepts:
																									Value  	Std. Error t value 
				Negative|No Change  -0.7312   0.0476   -15.3597
				No Change|Positive  -0.7105   0.0475   -14.9554
			\end{verbatim}
			
		\begin{verbatim}
								  REG OIL Negative No Change Positive
								1   0   0   0.3249    0.0046   0.6705
								2   1   0   0.2442    0.0038   0.7519
								3   0   1   0.3699    0.0048   0.6252
								4   1   1   0.2827    0.0042   0.7131
		\end{verbatim}
			
		The ordered model now works on a single latent scale for GDP difference such that it goes from Negative to No Change to Positive and each 'jump' on the scale is seen as the same. Unlike the unordered model, there is only one effect per variable, where a one-unit increase in $X$ is associated with on average the $\beta$ change in log-odds of being in a higher category in the order. This model relies on proportional odds assumption however that the effect of covariates is assumed to be the same assumed to be the same across the different categories.  \\
		
		\textbf{Cut-off Points}
		The cut-off point represents where the model (assuming a latent continuous variable underpinning outcome categories order) crosses certain thresholds where the 'boundaries' are between categories. The first threshold (-0.73) separates the 'Negative' GDP difference to the 'No Change' outcome, and the second (-0.71) is the threshold that distinguishes 'No Change' to 'Positive'. These are used when calculating predicted probabilities, such that the probability corresponding to covariate values assigns the categorical outcomes. \\
		
		\textbf{Coefficients}
		Holding all else constant, being a democracy is associated with, on average, a 0.39 increases the log-odds of being in a higher GDP difference category, i.e. the odds of them being in a higher (more positive) GDP difference category is 50\% higher than non-democracies ($e^{0.398} \approx1.49$) . The effect is statistically signifiant, indicating democracies are more likely to experience more positive GDP outcomes (on ordered scale) and less likely to be in the negative GDP outcome categories. \\
		
		Oil on the other hand has a negative coefficient, suggesting it has a tendency to being associated with more negative outcomes but this effect is not significant at the $\alpha =0.05$ threshold therefore we cannot reliably distinguish its effect. 
		
		\textbf{Predicted Probabilities}
		Looking at the predicted probabilities, we can see similar to the unordered logit the outcomes are overwhelmingly Positive or Negative with No change in GDP remaining very unlikely across all scenarios. 
		
		The impact of democracy can be seen again comparing the first two rows and the last two rows to determine the change of regime holding oil constant. A country being a democracy consistently shifts the probability from negative to positive GDP outcomes, increasing the likelihood of having a positive GDP difference from previous year by about 8-9\%. 
		On the other hand, being a dominant oil exporter is associated with a shift towards more negative and fewer positive GDP differences, although to a smaller extent with the likelihood of having a negative change in GDP increasing by around 4-5\%. 
	
		
\end{enumerate}

\section*{Question 2} 
\vspace{.25cm}

\noindent Consider the data set \texttt{MexicoMuniData.csv}, which includes municipal-level information from Mexico. The outcome of interest is the number of times the winning PAN presidential candidate in 2006 (\texttt{PAN.visits.06}) visited a district leading up to the 2009 federal elections, which is a count. Our main predictor of interest is whether the district was highly contested, or whether it was not (the PAN or their opponents have electoral security) in the previous federal elections during 2000 (\texttt{competitive.district}), which is binary (1=close/swing district, 0="safe seat"). We also include \texttt{marginality.06} (a measure of poverty) and \texttt{PAN.governor.06} (a dummy for whether the state has a PAN-affiliated governor) as additional control variables. 

\begin{enumerate}
	\item [(a)]
	Run a Poisson regression because the outcome is a count variable. Is there evidence that PAN presidential candidates visit swing districts more? Provide a test statistic and p-value.
		
		\lstinputlisting[language=R, firstline=114, lastline= 126]{PS03_EB.R}	
		% Table created by stargazer v.5.2.3 by Marek Hlavac, Social Policy Institute. E-mail: marek.hlavac at gmail.com
		% Date and time: Sun, Mar 15, 2026 - 15:07:33
		\begin{table}[!htbp] \centering 
			\caption{Poisson Regression on the Number of times PAN visited a district} 
			\label{} 
			\begin{tabular}{@{\extracolsep{5pt}}lc} 
				\\[-1.8ex]\hline 
				\hline \\[-1.8ex] 
				& \multicolumn{1}{c}{\textit{Dependent variable:}} \\ 
				\cline{2-2} 
				\\[-1.8ex] & Number of PAN Visits \\ 
				\hline \\[-1.8ex] 
				Competitive District & $-$0.081 \\ 
				& (0.171) \\ 
				& \\ 
				Marginality & $-$2.080$^{***}$ \\ 
				& (0.117) \\ 
				& \\ 
				PAN-Affiliated Governor & $-$0.312$^{*}$ \\ 
				& (0.167) \\ 
				& \\ 
				Constant & $-$3.810$^{***}$ \\ 
				& (0.222) \\ 
				& \\ 
				\hline \\[-1.8ex] 
				Observations & 2,407 \\ 
				Log Likelihood & $-$645.606 \\ 
				Akaike Inf. Crit. & 1,299.213 \\ 
				\hline 
				\hline \\[-1.8ex] 
				\textit{Note:}  & \multicolumn{1}{r}{$^{*}$p$<$0.1; $^{**}$p$<$0.05; $^{***}$p$<$0.01} \\ 
			\end{tabular} 
		\end{table} 
		\begin{verbatim}
			Overdispersion test
			
			data:  poisson_model
			z = 1.0668, p-value = 0.143
			alternative hypothesis: true dispersion is greater than 1
			sample estimates:
			dispersion 
			2.09834  
		\end{verbatim}
			
		\vspace{1cm}
		
		\lstinputlisting[language=R, firstline=129, lastline= 131]{PS03_EB.R}	
		\begin{verbatim}
			z test of coefficients:
			
																									Estimate Std. Error  z value Pr(>|z|)    
			(Intercept)          -3.810235   0.222095 -17.1559  < 2e-16 ***
			competitive.district -0.081352   0.170688  -0.4766  0.63364    
			marginality.06       -2.080144   0.117339 -17.7277  < 2e-16 ***
			PAN.governor.06      -0.311579   0.166731  -1.8688  0.06166 .  
			---
			Signif. codes:  0 ‘***’ 0.001 ‘**’ 0.01 ‘*’ 0.05 ‘.’ 0.1 ‘ ’ 1
		\end{verbatim}
		
		By estimating a Poisson regression we can see the relationship between districts characteristics and the number of visits made it it by the Mexican presidential candidate. The results of dispersion test ($p =0.143$) does not provide sufficient evidence to reject the null meaning we can conclude the data is not overdispersed (i.e. the variance is approximately the mean) and that a regular Poisson is appropriate. 
		
		There is not sufficient evidence to conclude that PAN presidential candidates visit swing districts more (competitive.district = 1). Holding all other variables constant, competitive districts were not significantly associated with the number of visits ($\beta = -0.081, z = -0.476, p = 0.634$). Since $p>0.05$ we cannot reject the null that competitive districts have no effect on the number of presidential PAN candidate visits.  
		 
	\item [(b)]
	Interpret the \texttt{marginality.06} and \texttt{PAN.governor.06} coefficients. \\
	
	Holding all other variables constant, a one unit increase in Marginality is associated with on average an multiplicative decrease of the log count of PAN visits of -2.08, in other words a roughly 87.5\% decrease in the expected number of visits ($e^{-2.08} \approx .125, 1-.125 = 87.5\%$). There is sufficient evidence therefore to reject the null, and conclude that the measure of poverty in a district is significantly associated with the number of PAN visits by the Mexican candidate ($ z = -17.73, p < 2e-16$).
	
	Holding all other variables constant, having a PAN-affiliated Governor in that district is associated with, on average, a multiplicative decrease of the log count of PAN visits by -0.312, i.e. a roughly 27\% roughly lower expected count ($e^{-.312}\approx0.73$). This effect however is not statistically significant at the standard $\alpha = 0.05$ level ($z = -1.87, p = .061$), indicating weak evidence to suggest that a district having a PAN affiliated governor impacts the number of candidate visits. 
	
	\item [(c)]
	Provide the estimated mean number of visits from the winning PAN presidential candidate for a hypothetical district that was competitive (\texttt{competitive.district}=1), had an average poverty level (\texttt{marginality.06} = 0), and a PAN governor (\texttt{PAN.governor.06}=1). \\
	
	\lstinputlisting[language=R, firstline=133, lastline= 138]{PS03_EB.R}	
	
	The model predicts that for the winning PAN candidate from a hypothetical swing district with average poverty level and a PAN-affiliated governor the expected mean number of visits would be 0.0149. This indicates a very low expected occurrence which makes sense given the high frequency of districts with 0 visits.  The expected mean $\lambda = 0.0149$ indicates that the probability for this district to have no visits by the candidate is 98.5\% ($e^{-.0149} \approx .986$) and only around a 1.5\% chance of 1 or more visits by the candidate (given the district characteristics).
	
	\begin{Verbatim}
		   1 
		0.0149
	\end{Verbatim}
	
\end{enumerate}

\end{document}
